\section{Implications éthiques, déontologiques et légales de l’incident }

La problématique qui m’interpelle d’emblée est le manque de prise en charge lors de la naissance de l’enfant. Suivre le protocole de naissance classique ne me semble pas adapté à la situation. La physiopathologie du déni de grossesse n’est pas encore élucidée entièrement à ce jour mais il est bien connu qu’il est le reflet d’un traumatisme chez la maman. De plus, le déni de grossesse total, c’est-à-dire jusqu’à l’accouchement, représente un facteur de risque du passage à l’acte et la maltraitance de l’enfant. L’accouchement est un bouleversement psychologique intense pour la maman. Elle se présente aux urgences, dans ce cas pour des lombalgies puis se retrouve engagée dans un travail  en vue d’une naissance. L’enfant n’est pas attendu, aucune relation n’a été créée lors de la grossesse, la femme n’a pas eu l’occasion de devenir mère comme l’explique Sophie Marinopoulos dans \textcite{marinopoulos_deni_2019}. A la naissance du patient, il y a un manquement clair au niveau de plusieurs principes de la Bioéthique selon Beauchamp et Childress. D’une part la bienfaisance vis-à-vis de l’enfant a été négligée vu les facteurs de risque de négligence et la prévention d’un point de vue médical qui était possible. D’autre part, la bienfaisance vis-à-vis de la mère. Le déni de grossesse montrant clairement une souffrance psychique qui aurait dû être explorée par un professionnel de santé formé.\\  

Une revue de la littérature,  \textcite{lagedamon2025accompagnement} appuie une fois de plus l’intérêt de prendre en charge la maman qui a souffert d’un déni de grossesse. Un protocole standardisé afin de prendre en charge le déni de grossesse manque cruellement ce qui a pour conséquence des prises en charge hétérogènes voire inexistantes. Or, un suivi psychologique pour la maman qui a un risque augmenté de stress post-traumatique ou dépression post partum, a montré des bénéfices. Les psychothérapies telles que EMDR, \textit{Eye Movement Desensitization and Reprocessing}, peuvent permettre une amélioration clinique de la maman. Ensuite, l’établissement d’un lien d’attachement avec son enfant est envisageable et peut permettre de prévenir les situations de négligence voire les néonaticides qui sont certes rares mais extrêmement graves.\\

Je remets également en question le respect du principe de la justice dans ce contexte. Autrement appelé l’équité, la justice consiste à traiter chaque personne de façon loyale, équitable et appropriée en fonction de ses spécificités. L’enfant a le droit d’avoir les mêmes chances de s’épanouir malgré un contexte de déni de grossesse et un milieu défavorisé. Est-il réellement éthique de ne pas avoir mis en place des aides autours de la famille afin de compenser l’environnement de l’enfant au moment de sa naissance ?\\

Je pourrais également ajouter à ce principe de justice comment peut-on demander à des parents ayant vécu un déni de grossesse total d’élever leur enfant comme une naissance classique sans outils, sans aide extérieure ?\\

Ce cas clinique met en lumière toute la complexité du signalement pour la protection de l’enfant dans un cadre de maltraitance. Le placement d’enfant soulève de nombreux défis éthiques. Le conflit entre la bienfaisance et la non-malfaisance est au centre du dilemme éthique. D’une part, la priorité est le bien-être de l’enfant voire la survie de celui-ci. Son passage en réanimation témoigne de l’urgence de la situation. La cassure de courbe de ce patient a été le point d’appel pour établir le diagnostic de malnutrition par manque d’apport. En effet, après la mise en évidence des facteurs de risque prédisposant à la maltraitance tel que l’état psychologique de la mère, le contexte de déni de grossesse, la situation financière de la famille, les connaissances des parents à propos du développement de l’enfant,…  les pédiatres responsables devaient évoquer la maltraitance comme expliqué dans ce clinical report \textcite{block2005failure}. De plus, la déontologie appuie d’autant plus le principe de bienfaisance. L’article 61 du code de Déontologie Médicale belge stipule que « \textit{Si la situation le justifie, et pour autant que la personne vulnérable capable de discernement y consente, le médecin s'adressera à un confrère compétent en la matière ou fera appel à une structure pluridisciplinaire spécifiquement établie pour gérer cette problématique. Si la personne vulnérable est menacée par un danger grave et imminent ou s'il y a des indices graves d'un danger sérieux et réel que d'autres personnes vulnérables soient victimes de maltraitance ou négligence et que le médecin n'a pas d'autre moyen d'offrir une protection, il peut avertir le procureur du Roi de ses constatations }». Suite à l’échec de la mise en place d’aide au domicile, une stratégie de non malfaisance dans le but de conserver la relation thérapeutique avec les parents, le signalement fut une solution envisageable pour l’ensemble de l’équipe médicale.\\

L’aspect légal intervient également dans la prise de décision du placement de l’enfant. La loi qui doit être prise en compte dans cette situation. Selon l'article $15, § 2$ de la loi du $22$ août $2002$ relative aux droits du patient, le médecin a le droit de déroger à la décision des parents (ou du représentant légal) si celle-ci met en danger la vie de l'enfant ou constitue une atteinte grave à sa santé. Donc, en considérant l’aspect légal, le corps médical a l’obligation de faire ce signalement dans l’immédiat afin d’assurer la survie du patient.\\ 

Toutefois, le principe de non-malfaisance doit être exploré dans cette situation. En effet, est-ce que le placement de l’enfant est la meilleure solution pour lui et son développement ? Il est démontré dans cette étude \textcite{poitras2017contacts} que le placement des jeunes enfants réduit la probabilité de réunification avec les parents biologiques et une stabilité de trajectoire. Selon eux, la stabilité de la trajectoire se définit par trois critères : les tentatives de réunification familiale, l’échec de la réunification et le nombre d’hébergements qu’a vécus l’enfant. Toutefois, plusieurs limites se retrouvent dans cette étude dont un échantillon de petite taille ce qui empêche de pouvoir interpréter les résultats avec certitude. Cependant, cette étude confirme également les résultats obtenus dans la littérature scientifique. Peut-on extrapoler cette instabilité de trajectoire comme représentation objective du malheur de l’enfant ou en tout cas d’une mauvaise prise en charge au final ?\\

Par contre, une autre étude, \textcite{mascaro2012evaluation} démontre l’importance d’apporter l’attention, la stimulation psychique et physique du nourrisson durant cette période critique. Sans cette attention, l’enfant peut développer un trouble important du développement d’où l’importance de sortir l’enfant de cet environnement négligent. La pouponnière peut être une option, si et seulement si elle peut assurer une prise en charge constante et régulière. Donc, la décision entre séparer l’enfant de sa famille et le placer dans une famille d’accueil ou une pouponnière est bien sûr à nuancer et comparer les soins qui sont possibles et à quelle fréquence concrètement ils peuvent être réalisés. Une autre remarque reprise dans cet article est l’importance de la persévérance. Celui-ci n’a effectué un suivi que sur 4 mois mais déjà dans ce délai, ils ont observé un rythme différent chez plusieurs enfants. Cela signifie donc, que le soignant responsable de l’enfant ne doit pas perdre espoir une fois les soins mis en place car des améliorations voire une reprise du développement physiologique est possible. Le dilemme éthique et la balance pour ou contre le placement de l’enfant n’est pas univoque, il s’agit d’un équilibre fragile.\\ 

En conclusion, le plus grand problème éthique que j’ai identifié dans cette situation clinique est le manque de prise en charge à la naissance du patient. Un autre aspect éthiquement intéressant est tout la complexité du placement de l’enfant pour sa survie et son bon développement ainsi que les alternatives de prises en charge possible pour l’enfant et le déni de grossesse de la maman.\\
